%=======================================================================
%  CHAPTER 6 — NUMERICAL PROBLEMS AND SOLUTIONS
%=======================================================================
\section*{\color{acc}Ch 6 \textperiodcentered{} Speech and Channel Coding \gap
\textcolor{sub}{Numerical Problems \& Solutions}}
\vspace{-4pt}{\color{acc}\rule{\textwidth}{1.1pt}}\vspace{4pt}

{\color{sub}\footnotesize The examined calculation here is the \textbf{7-bit Hamming code}
\yr{(\textbf{80 Ch} \m{5})}, worked both ways below. \textbf{Convolutional encoding, Viterbi
decoding} \tF{5} are asked as ``explain'', but they are only convincing with a worked trellis, so
one is given. The remaining problems are the tutorial-set coder-rate calculations, which are the
standard follow-ups to the speech coding theory.}

\lead{The formula sheet}
\begin{itemize}
  \item \textbf{Hamming code} \; $n = 2^m - 1$, \; $k = 2^m - m - 1$, \; $n-k = m$, \;
        $d_{\min} = 3$, \; $t = 1$, \; rate $= k/n$
  \item \textbf{Detect} $t$ errors: $d' \ge t+1$. \qquad \textbf{Correct} $t$ errors:
        $d' \ge 2t+1$
  \item \textbf{Convolutional code} \; rate $= k/n$, \; constraint length $K$, \;
        $2^{K-1}$ trellis states
  \item \textbf{Nyquist} \; a band-pass signal of bandwidth $B$ needs $f_s = 2B$ samples/s
  \item \textbf{Bit rate} \; $R = f_s \times$ (bits per sample)
  \item \textbf{Quantizer} \;
        $D = \displaystyle\int (x - f_Q(x))^2 p(x)\,dx$, \quad
        $S = E[x^2] = \displaystyle\int x^2 p(x)\,dx$, \quad
        $\text{SDR} = 10\log_{10}(S/D)$
  \item \textbf{Coding gain} \;
        $G = \left(\dfrac{E_b}{N_0}\right)_{\text{uncoded}}
           - \left(\dfrac{E_b}{N_0}\right)_{\text{coded}}$ dB, at the same BER
\end{itemize}

\hr

\T{Problem 1 \quad 7-bit Hamming code: encode and then detect}
{\color{sub}\footnotesize \tP{1} \yr{\textbf{80 Ch} \m{5}} \gap$\cdot$\gap
\mk{the only pure channel-coding calculation in 22 papers, so it is worth banking exactly}}

\begin{asked}{2080 Chaitra \textperiodcentered{} Q7 \textperiodcentered{} [3+5]}
Explain the characteristics of Speech. Explain 7-bit hamming code method with relevant example.
\end{asked}
{\color{sub}\footnotesize The paper prints no data word, so supply your own. The run below
encodes \textbf{1011} into a $(7,4)$ Hamming codeword, then receives a corrupted word and
locates the faulty bit --- that is the ``relevant example'' the 5 marks are for. \added}

\lead{Step 0: the structure, once}
\begin{itemize}
  \item For $m = 3$: \; $n = 2^3 - 1 = 7$, \; $k = 2^3 - 3 - 1 = 4$, \; parity bits $= 3$.
  \item Parity bits go in positions that are \textbf{powers of 2}: 1, 2, 4. Data fills the rest:
        3, 5, 6, 7.
  \item Layout, written with bit 7 on the \emph{left} and bit 1 on the \emph{right}:
        \[ D_7\; D_6\; D_5\; P_4\; D_3\; P_2\; P_1 \]
  \item Each parity bit checks the positions whose binary index contains its own bit:
\end{itemize}

\vspace{2pt}
\begin{center}\small
\begin{tabular}{@{}llll@{}}
\toprule
\hrow \thead{Parity bit} & \thead{Checks positions} & \thead{Binary indices} &
\thead{Why} \\
\midrule
$P_1$ & 1, 3, 5, 7 & 001, 011, 101, 111 & all have bit 0 set \\
$P_2$ & 2, 3, 6, 7 & 010, 011, 110, 111 & all have bit 1 set \\
$P_4$ & 4, 5, 6, 7 & 100, 101, 110, 111 & all have bit 2 set \\
\bottomrule
\end{tabular}
\end{center}

\vspace{2pt}
{\color{sub}\footnotesize \mk{This is the whole trick.} Because the groups are built on the binary
indices, the three parity checks read together \textbf{are} the binary address of the faulty bit.}

\lead{Part A \; Encode the data word 1011 with even parity}

\begin{itemize}
  \item \textbf{Step 1: place the data bits} into positions 7, 6, 5, 3:
        \[ D_7 = 1,\quad D_6 = 0,\quad D_5 = 1,\quad D_3 = 1 \]
  \item \textbf{Step 2: compute each parity bit} so that its group has an \textbf{even} number of
        ones.
\end{itemize}

\vspace{2pt}
\begin{center}\small
\begin{tabular}{@{}lllr@{}}
\toprule
\hrow \thead{Parity} & \thead{Group contents} & \thead{Ones already present} & \thead{Set to} \\
\midrule
$P_1$ & $(P_1,\,D_3,\,D_5,\,D_7) = (P_1,1,1,1)$ & 3 (odd) & $\mathbf{1}$ \\
$P_2$ & $(P_2,\,D_3,\,D_6,\,D_7) = (P_2,1,0,1)$ & 2 (even) & $\mathbf{0}$ \\
$P_4$ & $(P_4,\,D_5,\,D_6,\,D_7) = (P_4,1,0,1)$ & 2 (even) & $\mathbf{0}$ \\
\bottomrule
\end{tabular}
\end{center}

\begin{itemize}
  \item \textbf{Step 3: assemble the codeword}, bit 7 first:
\end{itemize}

\vspace{2pt}
\begin{center}\small
\begin{tabular}{@{}rccccccc@{}}
\toprule
\hrow Position & 7 & 6 & 5 & 4 & 3 & 2 & 1 \\
Contents & $D_7$ & $D_6$ & $D_5$ & $P_4$ & $D_3$ & $P_2$ & $P_1$ \\
\midrule
\textbf{Value} & \textbf{1} & \textbf{0} & \textbf{1} & \textbf{0} & \textbf{1} &
\textbf{0} & \textbf{1} \\
\bottomrule
\end{tabular}
\end{center}

\[ \boxed{\text{transmitted codeword} = 1\,0\,1\,0\,1\,0\,1} \]

\lead{Part B \; A codeword $1011011$ is received. Is it correct? If not, locate the error.}

\begin{itemize}
  \item \textbf{Step 1: lay out the received bits.}
\end{itemize}

\vspace{2pt}
\begin{center}\small
\begin{tabular}{@{}rccccccc@{}}
\toprule
\hrow Position & 7 & 6 & 5 & 4 & 3 & 2 & 1 \\
\textbf{Received} & \textbf{1} & \textbf{0} & \textbf{1} & \textbf{1} & \textbf{0} &
\textbf{1} & \textbf{1} \\
\bottomrule
\end{tabular}
\end{center}

\begin{itemize}
  \item \textbf{Step 2: recheck each parity group}, still assuming even parity.
\end{itemize}

\vspace{2pt}
\begin{center}\small
\begin{tabular}{@{}llllr@{}}
\toprule
\hrow \thead{Check} & \thead{Positions} & \thead{Bits read} & \thead{Parity} &
\thead{Syndrome bit} \\
\midrule
$P_1$ & 1, 3, 5, 7 & $1,0,1,1$ & \mk{ODD} $\rightarrow$ error & $\mathbf{1}$ \\
$P_2$ & 2, 3, 6, 7 & $1,0,0,1$ & EVEN $\rightarrow$ ok & $\mathbf{0}$ \\
$P_4$ & 4, 5, 6, 7 & $1,1,0,1$ & \mk{ODD} $\rightarrow$ error & $\mathbf{1}$ \\
\bottomrule
\end{tabular}
\end{center}

\begin{itemize}
  \item \textbf{Step 3: read the syndrome.} Write it as $P_4 P_2 P_1$:
        \[ \text{syndrome} = 1\,0\,1_{(2)} = 4 + 0 + 1 = 5 \]
  \item \textbf{Step 4: conclude.}
        \[ \boxed{\text{The received word is INCORRECT; bit position 5 is in error.}} \]
  \item \textbf{Step 5: correct it} by flipping bit 5, $1 \rightarrow 0$:
        \[ \text{corrected codeword} = 1\,0\,0\,1\,0\,1\,1 \]
        \[ \text{recovered data (positions 7,6,5,3)} = 1\,0\,0\,0 \]
  \item \mk{Non-zero syndrome $=$ error, and its decimal value is the position.
        A zero syndrome means no single-bit error.}
  \item \textbf{Limitation to state:} $d_{\min} = 3$, so this code corrects \textbf{one} bit and
        detects \textbf{two}. A \textbf{two-bit} error produces a non-zero syndrome pointing at
        the \emph{wrong} position, and ``correcting'' it makes things worse.
        \mk{That is why Hamming codes are always paired with interleaving} (Ch 5) on a fading
        channel, where errors arrive in bursts.
\end{itemize}

\lead{Hamming code parameters for other $m$}
\vspace{2pt}
\begin{center}\small
\begin{tabular}{@{}rrrrr@{}}
\toprule
\hrow $m$ & $n = 2^m-1$ & $k = 2^m-m-1$ & parity bits & rate $k/n$ \\
\midrule
3 & 7  & 4  & 3 & 0.571 \\
4 & 15 & 11 & 4 & 0.733 \\
5 & 31 & 26 & 5 & 0.839 \\
6 & 63 & 57 & 6 & 0.905 \\
\bottomrule
\end{tabular}
\end{center}
{\color{sub}\footnotesize \mk{The rate improves as $m$ grows}, but the code still corrects only
\textbf{one} bit per block, so the protection per bit gets \emph{weaker}. That trade-off is worth
one line in any Hamming answer.}

\hr

\T{Problem 2 \quad Convolutional encoding and Viterbi decoding}
{\color{sub}\footnotesize \tF{5} \yr{\textbf{74 Bh} \m{5}, \textbf{75 Bh}, 70 Ma,
\textbf{72 Ash} \m{3}, \textbf{73 Bh} \m{5}}}

\begin{asked}{5 papers, always as a short note}
Write short notes on: \textbf{Convolutional encoding and decoding}
\ \yr{(2074 Bhadra Q9 c \m{5})}
\par\vspace{1pt}
Write short notes on: \textbf{Viterbi decoding algorithm}
\ \yr{(2073 Bhadra Q9 c \m{5} \textperiodcentered{} 2072 Ashwin Q9 c \m{3}
\textperiodcentered{} 2075 Bhadra Q12 b \textperiodcentered{} 2070 Magh Q10 c)}
\end{asked}
\begin{asked}[PROBLEM]{not a PYQ \textperiodcentered{} the example that carries the short note \added}
For the rate 1/2, constraint length $K = 3$ convolutional encoder defined by
$X_1 = M + M_1 + M_2$ and $X_2 = M + M_2$ (mod-2 addition), encode the message \textbf{1011}.
Then decode the received sequence $Y = 11\ 01\ 10$ using the Viterbi algorithm, showing the
trellis and the surviving path metrics.
\end{asked}

\lead{Part A \; Encoding 1011}
\begin{itemize}
  \item State $=$ the contents of the two memory stages $(M_1 M_2)$, so there are
        $2^{K-1} = 4$ states: $a = 00$, $b = 01$, $c = 10$, $d = 11$.
  \item Two flush (tail) bits are appended so the encoder returns to state $a$.
\end{itemize}

\vspace{2pt}
\begin{center}\small
\begin{tabular}{@{}rccccc@{}}
\toprule
\hrow \thead{Input $M$} & \thead{State $(M_1M_2)$} & \thead{$X_1 = M{+}M_1{+}M_2$} &
\thead{$X_2 = M{+}M_2$} & \thead{Output} & \thead{Next state} \\
\midrule
1 & 00 & $1{+}0{+}0 = 1$ & $1{+}0 = 1$ & \textbf{11} & 10 \\
0 & 10 & $0{+}1{+}0 = 1$ & $0{+}0 = 0$ & \textbf{10} & 01 \\
1 & 01 & $1{+}0{+}1 = 0$ & $1{+}1 = 0$ & \textbf{00} & 10 \\
1 & 10 & $1{+}1{+}0 = 0$ & $1{+}0 = 1$ & \textbf{01} & 11 \\
0 & 11 & $0{+}1{+}1 = 0$ & $0{+}1 = 1$ & \textbf{01} & 01 \\
0 & 01 & $0{+}0{+}1 = 1$ & $0{+}1 = 1$ & \textbf{11} & 00 \\
\bottomrule
\end{tabular}
\end{center}

\[ \boxed{\text{encoded output} = 11\;10\;00\;01\;01\;11} \]

\begin{itemize}
  \item 4 message bits plus 2 flush bits give 6 input bits, and at rate 1/2 that is
        \textbf{12 output bits}. \checkmark
  \item \mk{Code rate including the tail} is $4/12 = 1/3$, not $1/2$. Real systems use long
        frames so the tail overhead is negligible. Worth a line.
\end{itemize}

\lead{Part B \; Viterbi decoding of $Y = 11\;01\;10$}
\begin{itemize}
  \item \textbf{Metric used:} \textbf{Hamming distance} between the branch's expected output and
        the received pair. Lower is better.
  \item At every node keep only the \textbf{smaller-metric} incoming path (the survivor) and
        discard the other.
\end{itemize}

\vspace{2pt}
\begin{center}\small
\begin{tabular}{@{}lllll@{}}
\toprule
\hrow \thead{Step} & \thead{Received} & \thead{State} & \thead{Path metric} &
\thead{Surviving input path} \\
\midrule
1 & 11 & 00 & 2 & 0 \\
  &    & 10 & \mk{0} & 1 \\
\midrule
2 & 01 & 00 & 3 & 0, 0 \\
  &    & 01 & 2 & 1, 0 \\
  &    & 10 & 3 & 0, 1 \\
  &    & 11 & \mk{0} & 1, 1 \\
\midrule
3 & 10 & 00 & 3 & 1, 0, 0 \\
  &    & 01 & 2 & 1, 1, 0 \\
  &    & 10 & 3 & 1, 0, 1 \\
  &    & 11 & \mk{0} & 1, 1, 1 \\
\bottomrule
\end{tabular}
\end{center}

\lead{Reading the table}
\begin{itemize}
  \item \textbf{Step 1.} From state $00$, input 0 gives output $00$ (distance 2 from the received
        $11$), input 1 gives output $11$ (distance 0). So state $10$ carries metric 0.
  \item \textbf{Step 2.} From state $10$, input 0 gives $10$ (distance 1 from $01$) and input 1
        gives $01$ (distance 0). The path $1,1$ into state $11$ keeps metric 0.
  \item \textbf{Step 3.} From state $11$, input 1 gives output $10$, which \textbf{matches the
        received $10$ exactly}, so the metric stays 0.
  \item \textbf{Survivor:} the path ending in state $11$ with total metric $\mathbf{0}$.
        \[ \boxed{\text{decoded message} = 1\,1\,1} \]
  \item \textbf{Verify by re-encoding} $1,1,1$ from state $00$:
        $11$, then $01$, then $10$ $\;=\; Y$ exactly. \checkmark
        \mk{Metric 0 means the received sequence was a valid codeword, ie no channel errors.}
  \item \mk{Why Viterbi is tractable:} after every step there are still only \textbf{4} survivors,
        never more, because each state keeps one path. An exhaustive tree search would have
        $2^{L}$ paths after $L$ bits.
\end{itemize}

\hr

\T{Problem 3 \quad Sub-band coder bit rate}
{\color{sub}\footnotesize \tP{1} \yr{\textbf{\texttt{80 Bh}} \m{6}} \gap$\cdot$\gap
the calculation behind ``sub-band coding with transmitter and receiver''}

\begin{asked}{2080 Bhadra \textperiodcentered{} Q6 (a) \textperiodcentered{} [6]}
Explain sub-band coding with transmitter and receiver.
\end{asked}
\begin{asked}[PROBLEM]{not a PYQ \textperiodcentered{} the rate sum behind the block diagram \added}
A sub-band coder splits the input into the four sub-bands tabulated below and samples each at
its band-pass Nyquist rate. Find the total encoding bit rate, and compare it with coding the
full band at the same resolution.
\end{asked}

\vspace{2pt}
\begin{center}\small
\begin{tabular}{@{}rrrrrr@{}}
\toprule
\hrow \thead{Sub-band} & \thead{Band (Hz)} & \thead{Bandwidth} & \thead{$f_s = 2B$} &
\thead{Bits/sample} & \thead{Rate (bps)} \\
\midrule
1 & 225--450   & 225 Hz & 450 & 4 & 1800 \\
2 & 450--900   & 450 Hz & 900 & 3 & 2700 \\
3 & 1000--1500 & 500 Hz & 1000 & 2 & 2000 \\
4 & 1800--2700 & 900 Hz & 1800 & 1 & 1800 \\
\midrule
\multicolumn{5}{r}{\textbf{Total}} & \textbf{8300} \\
\bottomrule
\end{tabular}
\end{center}

\begin{itemize}
  \item \textbf{Step 1: bandwidth of each band} $= f_{\text{high}} - f_{\text{low}}$.
  \item \textbf{Step 2: sampling rate.} For \textbf{perfect reconstruction} each band-pass signal
        must be sampled at \textbf{twice its own bandwidth}, not twice its highest frequency.
        \mk{That saving is the whole reason for the low-pass translation step in the sub-band
        coder.}
        \[ f_{s,1} = 2(450-225) = 450\ \text{samples/s}, \quad
           f_{s,2} = 2(900-450) = 900 \]
        \[ f_{s,3} = 2(1500-1000) = 1000, \quad
           f_{s,4} = 2(2700-1800) = 1800 \]
  \item \textbf{Step 3: total rate.}
        \[ R = 450(4) + 900(3) + 1000(2) + 1800(1) \]
        \[ = 1800 + 2700 + 2000 + 1800 = \boxed{8300\ \text{bps} = 8.3\ \text{kbps}} \]
  \item \mk{Note the bit allocation:} the \textbf{narrowest, lowest} band gets the \textbf{most}
        bits (4), and the widest, highest band gets the fewest (1). That follows directly from the
        \textbf{PSD property} of speech: low bands carry most of the energy and most of the
        perceptual weight.
  \item \textbf{Compare:} plain 8 kHz / 8-bit PCM of the same speech would need
        \textbf{64 kbps}. Sub-band coding achieves nearly \mk{8:1 compression} while remaining a
        waveform coder.
\end{itemize}

\hr

\T{Problem 4 \quad Upper bound on speech coder rate in an FDMA system}
{\color{sub}\footnotesize \yr{tutorial set; the calculation behind ``choosing a speech codec for
mobile communications''}}

\begin{asked}[PROBLEM]{not a PYQ \textperiodcentered{} the sum behind ``choosing a speech codec'' \added}
A digital mobile system uses 810--826 MHz for the forward channel and 940--956 MHz for the
reverse channel. 90\% of the bandwidth is given to traffic channels, and the system must support
at least 1150 simultaneous calls using FDMA. The modulation achieves a spectral efficiency of
1.68 bps/Hz, and channel impairments force rate 1/2 forward error correction. Find the upper
bound on the transmission bit rate that a speech coder in this system can use.
\end{asked}

\begin{itemize}
  \item \textbf{Step 1: total band.}
        \[ B_{\text{total}} = 826 - 810 = 16\ \text{MHz} \]
        {\color{sub}\footnotesize The reverse band is the same width, and FDMA pairs them, so work
        with one direction.}
  \item \textbf{Step 2: bandwidth actually available to traffic.}
        \[ B_{\text{traffic}} = 0.9 \times 16 = 14.4\ \text{MHz} \]
        \mk{The other 10\,\% goes to control channels and guard bands.}
  \item \textbf{Step 3: bandwidth per channel.} FDMA divides it among the simultaneous calls:
        \[ B_{\text{ch}} = \frac{14.4\ \text{MHz}}{1150} = 12.52\ \text{kHz} \]
        \[ \Rightarrow \boxed{\text{maximum channel bandwidth} = 12.5\ \text{kHz}} \]
  \item \textbf{Step 4: raw bit rate the channel can carry.}
        \[ R_{\text{raw}} = 1.68\ \frac{\text{bps}}{\text{Hz}} \times 12.52\ \text{kHz}
                          = 21.04\ \text{kbps} \]
  \item \textbf{Step 5: subtract the FEC overhead.} A rate 1/2 code doubles the bits, so only half
        the raw rate is available to the speech coder:
        \[ R_{\text{speech}} = 0.5 \times 21.04 = \boxed{10.52\ \text{kbps}} \]
  \item \mk{Interpretation:} the speech coder must run at \textbf{10.5 kbps or below}. A plain
        waveform coder cannot do that at toll quality, so this system needs a
        \textbf{hybrid or vocoder} class coder, eg RPE-LTP at 13 kbps (too fast here) or a CELP
        variant at 8 kbps (fits).
  \item \mk{This is exactly the reasoning that forces cellular systems onto low-rate coders},
        and it is the numerical answer to ``why do we need speech coding techniques''.
\end{itemize}

\hr

\T{Problem 5 \quad Quantizer mean square error and signal to distortion ratio}
{\color{sub}\footnotesize \yr{tutorial set; the calculation behind the \textbf{PDF} property of
speech \tS{7}}}

\begin{asked}[PROBLEM]{not a PYQ \textperiodcentered{} the PDF property of speech, made numerical \added}
The input signal to a quantizer has the probability density function $p(x) = x/32$ for
$0 \le x \le 8$ and zero elsewhere. The quantization levels are $\{1, 3, 5, 7\}$.
\begin{enumerate}[label=\alph*), leftmargin=1.6em]
  \item Compute the mean square error distortion at the quantizer output.
  \item Compute the output signal to distortion ratio.
  \item How would you redistribute the quantization levels to decrease the distortion, and for
        what input pdf would this quantizer already be optimal?
\end{enumerate}
\end{asked}

\begin{itemize}
  \item \textbf{Step 0: check the pdf is valid.}
        \[ \int_0^8 \frac{x}{32}\,dx = \left[\frac{x^2}{64}\right]_0^8 = \frac{64}{64} = 1
           \quad\checkmark \]
  \item \textbf{Step 1: set the decision boundaries.} With levels at 1, 3, 5, 7 the midpoints give
        boundaries at $\{0, 2, 4, 6, 8\}$, so the quantizer is \textbf{uniform} with step 2.
  \item \textbf{Step 2: mean square error distortion.}
        \[ D = \int_{-\infty}^{\infty}\left(x - f_Q(x)\right)^2 p(x)\,dx \]
        \[ = \int_0^2 (x-1)^2\frac{x}{32}dx + \int_2^4 (x-3)^2\frac{x}{32}dx
          + \int_4^6 (x-5)^2\frac{x}{32}dx + \int_6^8 (x-7)^2\frac{x}{32}dx \]
  \item \textbf{Step 3: evaluate one integral to show the method.} For the first,
        \[ \int_0^2 (x^2-2x+1)\frac{x}{32}dx
           = \frac{1}{32}\int_0^2 \left(x^3 - 2x^2 + x\right)dx \]
        \[ = \frac{1}{32}\left[\frac{x^4}{4} - \frac{2x^3}{3} + \frac{x^2}{2}\right]_0^2
           = \frac{1}{32}\left(4 - \frac{16}{3} + 2\right)
           = \frac{1}{32}\cdot\frac{2}{3} = \frac{1}{48} \]
        The other three evaluate to $\tfrac{3}{48}$, $\tfrac{5}{48}$ and $\tfrac{7}{48}$ by the
        same method.
  \item \textbf{Step 4: total distortion.}
        \[ D = \frac{1+3+5+7}{48} = \frac{16}{48} = \frac{1}{3}
             = \boxed{0.3333} \]
  \item \textbf{Step 5: signal power.}
        \[ S = E[x^2] = \int_0^8 x^2 \frac{x}{32}dx = \frac{1}{32}\left[\frac{x^4}{4}\right]_0^8
             = \frac{1}{32}\cdot\frac{4096}{4} = \boxed{32} \]
  \item \textbf{Step 6: signal to distortion ratio.}
        \[ \text{SDR} = 10\log_{10}\!\left(\frac{S}{D}\right)
                      = 10\log_{10}\!\left(\frac{32}{0.3333}\right)
                      = 10\log_{10}(96.0) = \boxed{19.82\ \text{dB}} \]
  \item \textbf{Step 7: how to reduce the distortion.}
  \begin{itemize}
    \item Note that the four sub-integrals came out as $1,3,5,7$ over 48:
          \mk{the top interval contributes seven times as much distortion as the bottom one}.
    \item That is because $p(x) = x/32$ \textbf{rises} with $x$, so most of the probability mass
          sits at large amplitudes.
    \item \textbf{Fix: use a non-uniform quantizer.} Place \textbf{more levels where the
          probability is high} (large $x$) and fewer where it is low (small $x$), ie shrink the
          step size towards $x = 8$.
    \item This is exactly the \textbf{PDF property} of speech being exploited, and it is the
          principle behind $\mu$-law and A-law companding.
  \end{itemize}
  \item \textbf{Step 8: for what pdf is this quantizer optimal?}
        A \textbf{uniform} quantizer is optimal when every interval contributes equally, which
        happens only when the input is \mk{uniformly distributed} over $[0,8]$, ie
        $p(x) = 1/8$.
\end{itemize}

\hr

\T{Problem 6 \quad Speech coder rates: what fits in what}
{\color{sub}\footnotesize \yr{supporting numbers for the vocoder \tS{10} and LPC \tS{8}
questions}}

{\color{sub}\footnotesize Reference figures. \mk{Quoting one correct rate is often worth a mark}
in an otherwise descriptive answer.}

\vspace{2pt}
\begin{center}\small
\begin{tabular}{@{}llll@{}}
\toprule
\hrow \thead{Coder} & \thead{Class} & \thead{Bit rate} & \thead{Note} \\
\midrule
PCM (toll reference)   & Waveform & 64 kbps  & 8 kHz, 8 bit, $\mu$-law \\
ADPCM                  & Waveform & 32 kbps  & G.726 \\
Sub-band coding        & Waveform & 9.6--32 kbps & 4 or 8 bands \\
Adaptive transform (ATC) & Waveform & 9.6--20 kbps & DCT based \\
Voice-excited vocoder  & Hybrid   & 7.2--9.6 kbps & no pitch extraction needed \\
GSM RPE-LTP            & Hybrid   & \mk{13 kbps} & full rate GSM \\
CELP                   & Vocoder  & \mk{4.8 kbps} & codebook excitation \\
LPC                    & Vocoder  & 4.8 kbps (good quality), 2.4 kbps & all-pole model \\
\bottomrule
\end{tabular}
\end{center}

\begin{itemize}
  \item \textbf{Compression achieved} against the 64 kbps PCM reference:
        ADPCM $2\times$, sub-band $\approx 7\times$, GSM $\approx 5\times$,
        CELP $\approx 13\times$.
  \item \mk{The pattern to remember:} rate falls, complexity rises, and robustness falls.
        A 4.8 kbps CELP frame damaged by the channel is far more audible than a damaged
        64 kbps PCM sample. \textbf{That is why the lowest-rate coders need the strongest channel
        coding}, which is the link between the two halves of this chapter.
\end{itemize}

\hr

\T{Problem 7 \quad Coding gain}
{\color{sub}\footnotesize \pill{gdL}{gd}{syllabus, no PYQ yet} \gap$\cdot$\gap included because
it is the number that justifies every code in the chapter}

\begin{asked}[PROBLEM]{not a PYQ \textperiodcentered{} the number that justifies every code here \added}
An uncoded BPSK link needs $E_b/N_0 = 10.5$ dB to reach a BER of $10^{-5}$. With a rate 1/2,
$K = 7$ convolutional code and soft-decision Viterbi decoding, the same BER is reached at
$E_b/N_0 = 4.5$ dB. Find the coding gain in dB and the transmit power saving it represents.
\end{asked}

\begin{itemize}
  \item \textbf{Step 1: coding gain.}
        \[ G = \left(\frac{E_b}{N_0}\right)_{\text{uncoded}}
             - \left(\frac{E_b}{N_0}\right)_{\text{coded}}
             = 10.5 - 4.5 = \boxed{6\ \text{dB}} \]
  \item \textbf{Step 2: what 6 dB is worth in power.}
        \[ 10^{6/10} = 3.98 \approx 4 \]
        \mk{The transmitter can run at one quarter of the power for the same quality.}
  \item \textbf{Step 3: what it is worth in range}, at path loss exponent $n = 3.5$:
        \[ 10^{6/(10 \times 3.5)} = 10^{0.1714} = 1.48 \]
        ie \textbf{48\,\% more range}, or $1.48^2 = 2.2$ times the area from one site.
  \item \textbf{Step 4: what it costs.} The rate 1/2 code \textbf{doubles the transmitted bit
        rate}, so it needs twice the bandwidth (or halves the user data rate).
  \item \mk{TCM exists precisely to escape that cost}: by combining coding with a larger signal
        constellation it delivers coding gains up to 6 dB with \textbf{no bandwidth expansion}.
        Turbo codes then push the gain further still, close to the Shannon bound.
\end{itemize}
